\documentclass{beamer}
\usepackage{amsmath}
\usepackage{amssymb}

% Choose a theme
\usetheme{Madrid}
\usecolortheme{default}
\usefonttheme{serif}

% --- CUSTOM FOOTER ---
% Redefine the footline to remove author/institute and keep title/page number
\makeatletter
\setbeamertemplate{footline}
{
  \leavevmode%
  \hbox{%
  \begin{beamercolorbox}[wd=.5\paperwidth,ht=2.25ex,dp=1ex,center]{title in head/foot}%
    \usebeamerfont{title in head/foot}\insertshorttitle
  \end{beamercolorbox}%
  \begin{beamercolorbox}[wd=.5\paperwidth,ht=2.25ex,dp=1ex,right]{date in head/foot}%
    \usebeamerfont{date in head/foot}\insertframenumber{} / \inserttotalframenumber\hspace*{2ex}
  \end{beamercolorbox}}%
  \vskip0pt%
}
\makeatother
% --- END CUSTOM FOOTER ---

\title[Lecture 18: Delay time tomography]{Lecture 18: Delay time tomography}
\author{David Al-Attar}
\institute{Michaelmas term 2024}
\date{}

\begin{document}

% --- INTRODUCTION ---

\begin{frame}
    \titlepage
\end{frame}

\begin{frame}{Outline and Motivation}
    We now extend ray theory to image 3-D variations in the Earth's structure.
    \begin{itemize}
        \item We will introduce \textbf{delay time tomography}, a technique to create quantitative images of lateral velocity variations inside the Earth.
        \pause
        \item Such tomographic models provide our primary observational constraint on processes like mantle convection.
        \pause
        \item This is a classic \textbf{linear inverse problem}, and we will discuss a practical method for solving it known as \textbf{regularised least squares}.
    \end{itemize}
\end{frame}

% --- DELAY TIMES ---

\begin{frame}{Delay Time Measurements}
    \begin{itemize}
        \item If the Earth were perfectly spherically symmetric, the travel time of a seismic wave would depend only on the distance and depth of the source.
        \pause
        \item In reality, we observe small deviations from this prediction due to lateral variations in velocity structure.
    \end{itemize}
    \pause
    \begin{alertblock}{Delay Time}
        We define the delay time, $ \delta T $, as the difference between the observed travel time ($T$) and the time predicted by a reference spherical Earth model ($T_0$).
        $$ \delta T = T - T_0 $$
        A positive $ \delta T $ means the wave arrived late, indicating it travelled through a region that is, on average, slower than the reference model. These delay times are our data.
    \end{alertblock}
\end{frame}

% --- LINEARISATION ---

\begin{frame}{Relating Delay Times to Velocity (1/2)}
    The relationship between travel time $T$ and velocity $v$ is non-linear. However, since the observed delay times are small, we can assume the Earth's true velocity ($v$) is only a small perturbation ($\delta v$) away from the reference velocity ($v_0$).
    $$ v = v_0 + \delta v $$
    \pause
    This allows us to linearise the problem.
    \pause
    \begin{block}{Fermat's Principle}
        Recall Fermat's Principle: the travel time of a ray between two points is stationary with respect to small perturbations of the ray path.
    \end{block}
\end{frame}

\begin{frame}{Relating Delay Times to Velocity (2/2)}
    \begin{itemize}
        \item When we perturb the velocity field by $ \delta v $, the ray path itself also changes.
        \pause
        \item The total change in travel time, $ \delta T $, comes from two effects: (1) the change in velocity along the path, and (2) the change in the path's geometry.
        \pause
        \item By Fermat's Principle, the first-order contribution from the change in path geometry is zero.
    \end{itemize}
    \pause
    \begin{alertblock}{The Linearised Relationship}
        Therefore, to first order, the delay time is simply the integral of the slowness perturbation ($ \delta(1/v) \approx -\delta v / v_0^2 $) along the \textbf{unperturbed ray path} in the reference model.
        $$ \delta T = -\int_{\text{ray}_0} \frac{\delta v}{v_0^2} ds $$
    \end{alertblock}
\end{frame}

% --- INVERSE PROBLEM SETUP ---

\begin{frame}{Setting Up the Inverse Problem (1/2)}
    Our goal is to find the continuous velocity perturbation function, $ \delta\alpha(\mathbf{x}) $, from a finite set of delay time measurements, $ \{\delta T_i\} $.
    \begin{itemize}
        \item This is a fundamentally \textbf{underdetermined} problem.
        \pause
        \item To make it solvable, we \textbf{parameterise} the model. We assume the velocity perturbation can be represented as a weighted sum of a finite number of known basis functions, $ \varphi_j(\mathbf{x}) $.
    \end{itemize}
    \begin{block}{Model Parameterisation}
        $$ \delta\alpha(\mathbf{x})=\sum_{j=1}^{m}\delta\alpha_{j}\varphi_{j}(\mathbf{x}) $$
        The unknowns are now the $m$ coefficients, $ \delta\alpha_j $.
    \end{block}
\end{frame}

\begin{frame}{Setting Up the Inverse Problem (2/2)}
    Substituting the parameterised model into our linearised relationship for each measurement $i$ gives:
    $$ \delta T_{i}=\sum_{j=1}^{m}\left(-\int_{\text{ray}_{i}}\frac{\varphi_{j}}{\alpha_{0}^{2}}ds\right)\delta\alpha_{j}+e_{i} $$
    \pause
    This is a system of linear equations.
    \begin{alertblock}{Matrix Formulation}
        We can write this compactly in matrix-vector form:
        $$ \mathbf{d} = \mathbf{A}\mathbf{m} + \mathbf{e} $$
        \begin{itemize}
            \item $ \mathbf{d} $ is the $n \times 1$ data vector of delay times, $ \delta T_i $.
            \item $ \mathbf{m} $ is the $m \times 1$ model vector of unknown coefficients, $ \delta\alpha_j $.
            \item $ \mathbf{e} $ is the vector of measurement errors.
            \item $ \mathbf{A} $ is the $n \times m$ sensitivity matrix (or kernel), where $ A_{ij} = -\int_{\text{ray}_{i}}(\varphi_{j}/\alpha_{0}^{2})ds $.
        \end{itemize}
    \end{alertblock}
\end{frame}

% --- REGULARISED LEAST SQUARES ---

\begin{frame}{Non-Uniqueness and the Null Space}
    Even after parameterisation, the problem $ \mathbf{d} = \mathbf{A}\mathbf{m} $ is typically non-unique.
    \begin{itemize}
        \item This occurs if the matrix $ \mathbf{A} $ has a non-trivial \textbf{null space}.
        \pause
        \item A vector $ \mathbf{m}_0 $ is in the null space if $ \mathbf{A}\mathbf{m}_0 = \mathbf{0} $.
        \pause
        \item If $ \mathbf{m} $ is a solution, then $ \mathbf{m} + a\mathbf{m}_0 $ is also a solution for any scalar $a$, because $ \mathbf{A}(\mathbf{m} + a\mathbf{m}_0) = \mathbf{A}\mathbf{m} + a\mathbf{A}\mathbf{m}_0 = \mathbf{d} + \mathbf{0} = \mathbf{d} $.
    \end{itemize}
    \pause
    \begin{exampleblock}{Physical Meaning of the Null Space}
        The null space corresponds to structures in the Earth that our set of rays does not sample. For example, a small velocity anomaly in a region where no rays pass cannot be detected and is part of the null space.
    \end{exampleblock}
\end{frame}

\begin{frame}{Solving the Problem (1/2): Regularisation}
    To overcome non-uniqueness, we need to add more information or constraints.
    \begin{itemize}
        \item A standard least-squares approach minimises the data misfit:
        $$ J_{data}(\mathbf{m}) = (\mathbf{A}\mathbf{m}-\mathbf{d})^{T}\mathbf{C}^{-1}(\mathbf{A}\mathbf{m}-\mathbf{d}) $$
        This is not enough, as any model in the null space gives the same minimum misfit.
        \pause
        \item We use \textbf{regularisation} by adding a second term that penalises models we consider "undesirable" (e.g., models that are too large or too rough).
    \end{itemize}
    \begin{alertblock}{Regularised Misfit Function}
        We minimise a combined objective function:
        $$ J(\mathbf{m}) = \underbrace{(\mathbf{A}\mathbf{m}-\mathbf{d})^{T}\mathbf{C}^{-1}(\mathbf{A}\mathbf{m}-\mathbf{d})}_{\text{Data Misfit}} + \underbrace{\lambda \, \mathbf{m}^{T}\mathbf{B}\mathbf{m}}_{\text{Model Norm / Penalty}} $$
        Here, $ \lambda $ is a trade-off parameter.
    \end{alertblock}
\end{frame}

\begin{frame}{Solving the Problem (2/2): The Solution}
    To find the model $ \mathbf{m} $ that minimises the regularised misfit function $ J(\mathbf{m}) $, we set its gradient with respect to $ \mathbf{m} $ to zero.
    $$ \nabla J(\mathbf{m}) = \mathbf{0} $$
    \pause
    This yields a system of linear equations for $ \mathbf{m} $.
    \begin{block}{The Normal Equations}
        $$ (\mathbf{A}^{T}\mathbf{C}^{-1}\mathbf{A}+\lambda \mathbf{B})\mathbf{m} = \mathbf{A}^{T}\mathbf{C}^{-1}\mathbf{d} $$
    \end{block}
    \pause
    \begin{alertblock}{The Regularised Least-Squares Solution}
        Solving for $ \mathbf{m} $ gives the unique solution:
        $$ \mathbf{m}=(\mathbf{A}^{T}\mathbf{C}^{-1}\mathbf{A}+\lambda \mathbf{B})^{-1}\mathbf{A}^{T}\mathbf{C}^{-1}\mathbf{d} $$
        The regularisation term $ \lambda\mathbf{B} $ ensures the matrix is invertible.
    \end{alertblock}
\end{frame}

% --- BAYESIAN PERSPECTIVE ---

\begin{frame}{A Bayesian Perspective}
    An alternative way to think about inverse problems is through \textbf{Bayes' Theorem}.
    $$ p(\mathbf{m}|\mathbf{d}) \propto p(\mathbf{d}|\mathbf{m})p(\mathbf{m}) $$
    \pause
    \begin{itemize}
        \item $ p(\mathbf{m}) $ is the \textbf{prior}: what we believe about the model before seeing the data.
        \item $ p(\mathbf{d}|\mathbf{m}) $ is the \textbf{likelihood}: the probability of observing our data given a specific model.
        \item $ p(\mathbf{m}|\mathbf{d}) $ is the \textbf{posterior}: our updated knowledge of the model after considering the data.
    \end{itemize}
    \pause
    \begin{alertblock}{Connection to Least Squares}
        If we assume the prior knowledge and the data errors are both Gaussian, then the model that maximises the posterior probability is \textbf{identical} to the regularised least-squares solution.
        \begin{itemize}
            \item The regularisation term $ \lambda\mathbf{B} $ is equivalent to assuming a Gaussian prior on the model parameters.
        \end{itemize}
    \end{alertblock}
\end{frame}

% --- SUMMARY ---

\begin{frame}{Summary: What You Need to Know}
    \begin{itemize}
        \item \textbf{Delay Time Tomography} uses travel-time anomalies ($ \delta T = T_{obs} - T_{ref} $) to image 3-D velocity structure.
        \pause
        \item \textbf{Fermat's Principle} provides the linear relationship between delay time and velocity perturbations: $ \delta T = -\int (\delta v / v_0^2) ds $.
        \pause
        \item The inverse problem is set up as a linear system $ \mathbf{d} = \mathbf{A}\mathbf{m} + \mathbf{e} $ by parameterising the model.
        \pause
        \item This system is non-unique due to a \textbf{null space}. It is solved using \textbf{regularised least squares}, which balances data fit with a model penalty.
        \pause
        \item The regularised least-squares solution is equivalent to the most probable model in a \textbf{Bayesian} framework with Gaussian assumptions.
    \end{itemize}
\end{frame}

\end{document}
